\documentclass[aspectratio=169, table, xcdraw]{beamer}

\usetheme{Madrid}
\usecolortheme{default}
\definecolor{ucbblue}{RGB}{0, 51, 102}

\setbeamercolor{structure}{fg=ucbblue}
\setbeamercolor*{palette primary}{fg=white,bg=ucbblue}
\setbeamercolor*{palette secondary}{fg=white,bg=ucbblue!80}
\setbeamercolor*{palette tertiary}{fg=white,bg=ucbblue!90}
\setbeamercolor{title}{fg=white, bg=ucbblue}
\setbeamercolor{frametitle}{fg=white, bg=ucbblue}
\setbeamercolor{item}{fg=ucbblue}

\usepackage[T1]{fontenc}
\usepackage[utf8]{inputenc}
\usepackage{tgpagella}
\usefonttheme{serif}
\usepackage[spanish, es-noshorthands]{babel}
\usepackage{amsmath, amssymb, bm}
\usepackage{graphicx}
\usepackage[most]{tcolorbox}
\usepackage{tikz}
\usepackage[american]{circuitikz}
\usetikzlibrary{shapes.geometric, arrows.meta, positioning, calc}

\title[Thévenin y Norton]{Redes Equivalentes: Teoremas de Thévenin y Norton}
\subtitle{IMT-121 Circuitos Electrónicos I — Semana 6, Sesión 3}
\author[B. Quiroga Turdera]{M.Sc. Ing. Bernardo Quiroga Turdera}
\institute[UCB Tarija]{Universidad Católica Boliviana ``San Pablo'' — Sede Tarija}
\date{Semana 6}

\begin{document}

\begin{frame}
    \titlepage
\end{frame}

\begin{frame}{Recuperación Inicial}
    \begin{tcolorbox}[colback=ucbblue!5, colframe=ucbblue, title=Diagnóstico Rápido]
        \begin{enumerate}
            \item ¿Qué diferencia física y de cálculo existe entre un \textbf{cortocircuito} y un \textbf{circuito abierto}?
            \item ¿Qué sucede con una \textit{fuente de tensión independiente} al desactivarla para aplicar superposición?
            \item Si cambiamos \textbf{todo} el interior de una radio por un circuito moderno que produce exactamente el mismo sonido y consume la misma corriente, ¿lo notaría el usuario externo?
        \end{enumerate}
    \end{tcolorbox}
    \vspace{0.4cm}
    \textbf{El hilo conductor de hoy:} \\
    $\text{Red Original} \rightarrow \text{Puerto } a-b \rightarrow \text{Relación Externa } V-I \rightarrow \text{Equivalente Thévenin/Norton}$
\end{frame}

\begin{frame}{El Concepto de Puerto y Equivalencia Externa}
    Un \textbf{puerto} es el par de terminales mediante el cual una red interactúa con el mundo exterior (una carga $R_L$).
    
    \vspace{0.3cm}
    \begin{columns}
        \begin{column}{0.48\textwidth}
            \centering
            \begin{circuitikz}[scale=0.8, transform shape, thick]
                \draw[fill=ucbblue!10, draw=ucbblue, rounded corners] (-1,-1) rectangle (3,3);
                \node at (1,1) [align=center, text width=2.5cm] {\textbf{Red Lineal Compleja}};
                \draw (3,2) to[short, -o] (4,2) node[above]{$a$};
                \draw (3,0) to[short, -o] (4,0) node[below]{$b$};
                \draw (4,2) -- (5,2) to[R, l={$R_L$}] (5,0) -- (4,0);
                \draw[-Stealth, red] (5,2) -- (4.2,2) node[midway, above]{\small $I_L$};
                \draw[<->, blue] (4.2,1.8) -- (4.2,0.2) node[midway, left]{\small $V_{ab}$};
            \end{circuitikz}
        \end{column}
        \begin{column}{0.48\textwidth}
            \centering
            \begin{circuitikz}[scale=0.8, transform shape, thick]
                \draw[fill=green!10, draw=green!60!black, rounded corners] (0,-1) rectangle (3,3);
                \node at (1.5,1) [align=center, text width=2cm] {\textbf{Caja Ficticia}};
                \draw (3,2) to[short, -o] (4,2) node[above]{$a$};
                \draw (3,0) to[short, -o] (4,0) node[below]{$b$};
                \draw (4,2) -- (5,2) to[R, l={$R_L$}] (5,0) -- (4,0);
                \draw[-Stealth, red] (5,2) -- (4.2,2) node[midway, above]{\small $I_L$};
                \draw[<->, blue] (4.2,1.8) -- (4.2,0.2) node[midway, left]{\small $V_{ab}$};
            \end{circuitikz}
        \end{column}
    \end{columns}
    
    \vspace{0.4cm}
    \textbf{Condición de Equivalencia:} Dos redes son equivalentes si generan \textbf{exactamente la misma relación} entre $V_{ab}$ e $I_L$. A la carga $R_L$ no le importan las variables internas.
\end{frame}

\begin{frame}{Equivalente de Thévenin}
    Cualquier red resistiva lineal vista desde un puerto $a-b$ puede reemplazarse por una fuente ideal de tensión $\bm{V_{th}}$ en serie con una resistencia $\bm{R_{th}}$.

    \begin{columns}
        \begin{column}{0.48\textwidth}
            \centering
            \begin{circuitikz}[scale=0.8, transform shape, thick]
                \draw (0,0) to[V, l={$V_{th}$}] (0,2) to[R, l={$R_{th}$}] (2,2) to[short, -o] (3,2) node[above]{$a$};
                \draw (0,0) to[short, -o] (3,0) node[below]{$b$};
                \draw[<->, blue] (3,1.8) -- (3,0.2) node[midway, left]{\small $V_{oc} = V_{th}$};
            \end{circuitikz}
            \textbf{Paso 1: $\bm{V_{th}}$} \\
            Tensión de \textbf{circuito abierto} ($I_{\text{puerto}} = 0$).
        \end{column}
        \begin{column}{0.48\textwidth}
            \centering
            \begin{circuitikz}[scale=0.8, transform shape, thick]
                \draw (0,0) -- (0,2) node[midway, left]{\small Corto} to[R, l={$R_{th}$}] (2,2) to[short, -o] (3,2) node[above]{$a$};
                \draw (0,0) to[short, -o] (3,0) node[below]{$b$};
                \draw[Stealth-, red] (2.8,1.8) -- (4,1.8) node[right]{Mirada};
            \end{circuitikz}
            \textbf{Paso 2: $\bm{R_{th}}$} \\
            Resistencia equivalente desde $a-b$ con \textbf{fuentes independientes desactivadas}.
        \end{column}
    \end{columns}
\end{frame}

\begin{frame}{Equivalente de Norton}
    La \textbf{misma red} puede reemplazarse por una fuente ideal de corriente $\bm{I_N}$ en paralelo con una resistencia $\bm{R_N}$.
    
    \begin{columns}
        \begin{column}{0.48\textwidth}
            \centering
            \begin{circuitikz}[scale=0.8, transform shape, thick]
                \draw (0,0) to[I, l={$I_N$}] (0,2) -- (2,2) to[short, -o] (3,2) node[above]{$a$};
                \draw (1.5,2) to[R, l={$R_N$}] (1.5,0);
                \draw (0,0) to[short, -o] (3,0) node[below]{$b$};
                \draw[red] (3,2) -- (3.5,2) -- (3.5,0) -- (3,0);
                \node[red, right] at (3.5,1) {$I_{sc} = I_N$};
            \end{circuitikz}
        \end{column}
        \begin{column}{0.48\textwidth}
            \begin{itemize}
                \item $\bm{I_N}$: Corriente de \textbf{cortocircuito} del puerto ($V_{ab} = 0$).
                \item $\bm{R_N}$: Es exactamente idéntica a $\bm{R_{th}}$.
            \end{itemize}
            \vspace{0.3cm}
            \begin{tcolorbox}[colback=white, colframe=ucbblue, title=Relación Thévenin $\leftrightarrow$ Norton]
                \centering
                $R_N = R_{th}$ \\[0.2cm]
                $V_{th} = I_N \cdot R_{th}$
            \end{tcolorbox}
        \end{column}
    \end{columns}
\end{frame}

\begin{frame}{Ejemplo Guiado: Red y Planteamiento}
    Hallar equivalentes de Thévenin y Norton respecto a la carga $R_L$ (terminales $a-b$).
    
    \begin{center}
    \begin{circuitikz}[scale=0.85, transform shape, thick]
        % Red interna
        \draw (0,0) to[V, l={$12\,\text{V}$}] (0,2) to[R, l={$4\,\text{k}\Omega$}] (3,2);
        \draw (2.5,2) to[short, *-o] (6,2) node[above]{$a$};
        \draw (2.5,2) to[R, l_={$8\,\text{k}\Omega$}] (2.5,0);
        \draw (4.5,0) to[I, l={$1\,\text{mA}$}] (4.5,2);
        
        \draw (0,0) to[short, -o] (6,0) node[below]{$b$};
        \draw (2.5,0) node[ground]{};
        
        % Carga (removible)
        \draw[dashed, blue] (5,2) -- (6,2) to[R, l={$R_L$}] (6,0) -- (5,0);
    \end{circuitikz}
    \end{center}
    \textbf{Primer paso obligatorio:} Retirar la carga $R_L$ de los terminales $a-b$. Dejar el puerto abierto.
\end{frame}

\begin{frame}{Ejemplo Guiado: Paso 1 — Hallar $V_{th}$}
    Buscamos la tensión en los terminales abiertos $a-b$. Como no fluye corriente hacia el puerto, $V_{th}$ es simplemente la tensión del nodo interno superior respecto a tierra.
    
    \begin{columns}
        \begin{column}{0.4\textwidth}
            \centering
            \begin{circuitikz}[scale=0.75, transform shape, thick]
                \draw (0,0) to[V, l={$12\,\text{V}$}] (0,2) to[R, l={$4\,\text{k}\Omega$}] (3,2);
                \draw (2.5,2) to[short, *-o] (6,2) node[above]{$a$};
                \draw (2.5,2) to[R, l_={$8\,\text{k}\Omega$}] (2.5,0);
                \draw (4.5,0) to[I, l={$1\,\text{mA}$}] (4.5,2);
                \draw (0,0) to[short, -o] (6,0) node[below]{$b$};
                \draw (2.5,0) node[ground]{};
                \node at (5.5,1) {\textcolor{blue}{$V_{th}$}};
            \end{circuitikz}
        \end{column}
        \begin{column}{0.55\textwidth}
            \textbf{Opción 1: KCL en el nodo $a$:}
            \begin{equation*}
                \frac{V_a - 12}{4\,\text{k}} + \frac{V_a}{8\,\text{k}} - 1\,\text{mA} = 0 \implies V_a = \mathbf{10.667\,\text{V}}
            \end{equation*}
            \textbf{Opción 2: ¡Superposición!}
            \begin{itemize}
                \item Solo 12V ($I$ abierto): $V_a = 12 \cdot \frac{8}{12} = 8\,\text{V}$.
                \item Solo 1mA ($V$ corto): $V_a = 1\,\text{mA} \cdot (4\,\text{k} \parallel 8\,\text{k}) = 2.667\,\text{V}$.
                \item Total: $V_{th} = \mathbf{10.667\,\text{V}}$.
            \end{itemize}
        \end{column}
    \end{columns}
\end{frame}

\begin{frame}{Ejemplo Guiado: Paso 2 — Hallar $R_{th}$}
    Desactivamos TODAS las fuentes independientes y "miramos" la resistencia equivalente desde los terminales $a-b$ hacia adentro.
    
    \begin{columns}
        \begin{column}{0.4\textwidth}
            \centering
            \begin{circuitikz}[scale=0.75, transform shape, thick]
                \draw (0,0) -- (0,2) node[midway, left]{\small Corto} to[R, l={$4\,\text{k}\Omega$}] (3,2);
                \draw (3,2) to[short, *-o] (4.5,2) node[above]{$a$};
                \draw (3,2) to[R, l={$8\,\text{k}\Omega$}] (3,0);
                \node at (1.5,1) {\small Abierto};
                \draw (0,0) to[short, -o] (4.5,0) node[below]{$b$};
                \draw (2.5,0) node[ground]{};
                \draw[Stealth-, red] (4.2,1) -- (5.5,1) node[right]{$R_{th}$};
            \end{circuitikz}
        \end{column}
        \begin{column}{0.55\textwidth}
            \begin{itemize}
                \item La fuente de 12V se vuelve un cable (Corto).
                \item La fuente de 1mA se vuelve una interrupción (Abierto).
            \end{itemize}
            Desde $a-b$, vemos claramente que la resistencia de $4\,\text{k}\Omega$ y la de $8\,\text{k}\Omega$ están en paralelo.
            \begin{equation*}
                R_{th} = 4\,\text{k}\Omega \parallel 8\,\text{k}\Omega = \frac{4 \cdot 8}{4 + 8} = \mathbf{2.667\,\text{k}\Omega}
            \end{equation*}
        \end{column}
    \end{columns}
\end{frame}

\begin{frame}{Ejemplo Guiado: Paso 3 — Equivalentes y Verificación con Carga}
    \begin{columns}
        \begin{column}{0.45\textwidth}
            \textbf{Thévenin} \\
            $V_{th} = 10.667\,\text{V}$, $R_{th} = 2.667\,\text{k}\Omega$.
            
            \vspace{0.2cm}
            \textbf{Norton} \\
            $R_N = 2.667\,\text{k}\Omega$. \\
            $I_N = \frac{V_{th}}{R_{th}} = \frac{10.667}{2.667\,\text{k}} = \mathbf{4\,\text{mA}}$.
        \end{column}
        \begin{column}{0.5\textwidth}
            \begin{block}{Prueba Definitiva: Reconectando $R_L = 8\,\text{k}\Omega$}
                Si conectamos $8\,\text{k}\Omega$ al circuito Thévenin:
                \begin{equation*}
                    I_L = \frac{10.667\,\text{V}}{2.667\,\text{k} + 8\,\text{k}} = \mathbf{1\,\text{mA}}
                \end{equation*}
                \begin{equation*}
                    V_L = 1\,\text{mA} \cdot 8\,\text{k}\Omega = \mathbf{8\,\text{V}}
                \end{equation*}
                \textit{(Puedes comprobar resolviendo la red original entera con los $8\,\text{k}\Omega$ conectados y dará exactamente lo mismo).}
            \end{block}
        \end{column}
    \end{columns}
\end{frame}

\end{document}
